\section{Introduction}\label{sec:introduction}

A proper $k$-coloring of the Euclidean plane is a function
$c:\R^2\to\{1,\ldots,k\}$ such that $c(x)\ne c(y)$ whenever
$\norm{x-y}=1$. The Hadwiger--Nelson problem asks for the least such
$k$, denoted by $\chi(\R^2)$, with no regularity imposed on the color
classes. Equivalently, it asks for the chromatic number of the graph
whose vertices are all points of the plane and whose edges are the
unit-distance pairs. Edges in a finite drawing of this graph may cross.

The problem dates to Edward Nelson in 1950. Nelson obtained the lower
bound four, and John Isbell found the upper bound seven by a hexagonal
coloring; the historical record of these unpublished observations is
recounted by \citet[pp.~15--19]{Soifer2003History}. Hadwiger's published
formulation \citep{Hadwiger1961} asks for a covering by sets that avoid
unit distance and gives the bounds four and seven. The seven-point
construction of \citet{MoserMoser1961}, now called the Moser spindle,
gives a particularly small finite obstruction to three colors. We use
its two overlapping-triangle constraints again in the final step of our
proof and verify the required placement explicitly.

The lower bound remained four until \citet{deGrey2018} constructed a
finite unit-distance graph that is not four-colorable in 2018.
His corrected construction has 1,581 vertices. Subsequent work made
these finite obstructions smaller and their verification more accessible:
\citet{Heule2018} used satisfiability proofs and extraction of
unsatisfiable cores to obtain 553-vertex examples, and
\citet{Parts2020Minimization} obtained a 509-vertex graph.
\citet{ExooIsmailescu2020} gave an alternative proof of the lower bound five,
and Parts gave a human-verifiable proof \citep{Parts2020Human}.
These advances concern finite obstructions to four colors. The
compactness theorem of \citet[Theorem~1]{deBruijnErdos1951} says that a
graph is $k$-colorable, for fixed finite $k$, if and only if every finite
subgraph is $k$-colorable. Consequently an unrestricted lower bound can
also be established without first displaying a finite graph witnessing
it.

The problem with measurable color classes has a different history.
Write $\chi_m(\R^2)$ for the least number of colors when each class must
be Lebesgue measurable. \citet{Falconer1981} proved
$\chi_m(\R^2)\ge5$ using density points and rotations.
\citet{Payne2009} developed this method for translation-invariant
unit-distance subgraphs. In particular, the graph on all of $\R^2$
whose edges have rational unit displacement vectors is two-colorable,
while every measurable coloring requires at least five colors. Thus
measurable and unrestricted colorings cannot be identified merely
because their distance constraints look the same.

For colorings by regular regions, \citet{Woodall1973} studied
distance-avoiding closed covers and region decompositions. Townsend
corrected Woodall's argument for a six-color lower bound in the
planar-map setting, announcing the repair in 1981 and giving its details
in 2005 \citep{Townsend1981,Townsend2005}. More recently,
\citet{SokolovVoronov2025} proved a lower bound seven for polygonal
colorings in a locally finite map framework. Their interface argument
excludes both adjacent colors wherever distances to the interface take
values both below and above one. We derive corresponding exclusions
from weak measurable data, without assuming regular color boundaries.

We prove the following result, working throughout in ZFC.

\begin{theorem}\label{thm:main}
The Euclidean plane has no proper five-coloring, even when arbitrary
color classes are allowed. Consequently,
\[
  6\le \chi(\R^2)\le 7.
\]
\end{theorem}

The remaining alternatives six and seven are unresolved.\footnote{
The May 14, 2026 manuscript of \citet{Reed2026} announces
$\chi(\R^2)=7$ through a circle-density argument. Its proposed strict
bound of $\pi/3$ on the angular measure of a unit-independent subset
of the unit circle fails for the half-open arc
$\{e^{it}:0\le t<\pi/3\}$: this arc has angular measure $\pi/3$
